\chapter{Introducción}\label{sec:introduccion}

La clasificación de texturas depende de señales locales, distribuciones espaciales y respuestas en frecuencia que cambian con el dominio, la escala, la iluminación y el punto de vista \citep{cimpoi2014dtd,cimpoi2016fisher,liu2019survey}. Ninguna representación captura necesariamente todas esas variaciones. Los descriptores clásicos resumen propiedades diseñadas, como microestructuras, coocurrencias y gradientes \citep{ojala2002lbp,haralick1979glcm,dalal2005hog}; las redes convolucionales, los Transformers y los modelos autosupervisados producen características aprendidas con arquitecturas y objetivos de preentrenamiento diferentes \citep{he2016resnet,tan2019efficientnet,dosovitskiy2021vit,liu2021swin,oquab2024dinov2}. Esta diversidad crea una oportunidad para combinar representaciones, pero también introduce redundancia, escalas incompatibles y un aumento considerable de dimensionalidad.

La concatenación conserva la información de varios extractores y delega al clasificador la decisión sobre sus coordenadas. Agregaciones locales, Fisher vectors, bancos de filtros profundos y capas de codificación aprendida han mostrado que combinar respuestas puede mejorar el reconocimiento de texturas \citep{sanchez2013image,cimpoi2016fisher,zhang2017deepten}. Sin embargo, una mejora de la representación conjunta no demuestra que todos sus bloques sean necesarios ni que la complementariedad explique por sí sola el resultado. El mismo efecto puede aparecer al aumentar la dimensión o al reunir descriptores que ya son fuertes de manera individual.

Las arquitecturas de fusión y las combinaciones de representaciones motivan distinguir tres factores en la evaluación: calidad individual, cantidad de variables y composición del subconjunto. La distinción adquiere especial importancia en una biblioteca que reúne descriptores clásicos, CNN, Transformers y representaciones autosupervisadas o multimodales. La pregunta de estudio es \emph{cómo cambian el desempeño externo y la dimensionalidad al concatenar descriptores, frente a seleccionar uno individual, y qué aporta elegir la composición de esa concatenación}.

Para responder esta pregunta se evaluaron 21 descriptores sobre DTD, FMD, CUReT y Outex con SVM lineal y ResMLP. El mismo protocolo anidado compara cinco estrategias: Individual, Completa, Top-$k$, GFS y Homogénea. La comparación considera macro-F1 y exactitud, junto con la dimensión de las representaciones. Se completaron 28 condiciones externas por clasificador; cada una incluye las cinco estrategias sobre la misma biblioteca de 21 descriptores.

El trabajo aporta una comparación reproducible de estrategias de concatenación bajo particiones comunes, una evaluación no paramétrica entre datasets y una comparación del desempeño y la dimensionalidad de los subconjuntos seleccionados. GFS y Top-$k$ son procedimientos para elegir qué bloques concatenar dentro de esta comparación. El propósito es caracterizar las ganancias y los costos observados de combinar bloques, sin introducir un nuevo algoritmo de selección. La evidencia se interpreta atendiendo tanto a las diferencias observadas como a la incertidumbre estadística.

La tesis se organiza en seis capítulos: introducción, antecedentes, metodología, resultados, discusión y conclusión. El estudio anidado con cuatro datasets constituye la referencia actual. Los experimentos exploratorios previos se conservan en los materiales del proyecto y no se mezclan con sus contrastes estadísticos.

% ============================================================
% 2. TRABAJOS RELACIONADOS
% ============================================================
